\chapter{数据管理与数据库}

\section{领域导言与演进背景}

\begin{preamblebox}[知识领域概述与时代演进]
计算机科学的各个细分领域都可以精炼地概括为“研究面向特定目标的算法与数据结构”。对于数据管理 (DM)领域而言，这一目标正是：“\textbf{高效处理与组织持久化数据集——且这些数据集的规模往往过于庞大，无法完全容纳于计算机主存储器（内存）之中}”。

自 20 世纪 70 年代中期以来，本领域的专业教学与学术研究几乎排他性地聚焦于关系数据库系统。根据各高校不同的办学定位与培养侧重，以往学生在不同比例上系统修读了以下四个经典支柱：
\begin{enumerate}
    \item \textbf{数据建模与数据库设计}：例如实体-联系数据模型（E-R 模型）、关系模型、函数依赖与模式规范化范式理论。
    \item \textbf{声明式查询构建}：例如形式化关系代数、SQL 结构化查询语言。
    \item \textbf{查询处理与执行优化}：例如底层物理索引结构（B+ 树、静态与动态哈希）、核心查询物理执行算法（外部归并排序、选择、投影、多种表连接算法）、基于代价的查询优化器（等价代数变换、物理访问路径与索引选型）。
    \item \textbf{数据库管理系统（DBMS）内核实现}：例如并发控制与两阶段封锁协议、事务管理与 ACID 特性、内存缓冲池置换与脏页调度管理。
\end{enumerate}

\textbf{现代计算机专业人才培养定位的深刻变革：}\\
今天的行业与学术界对绝大多数计算机专业毕业生的核心期待，已明显转向要求其具备扎实的“\textbf{DBMS 高级用户与数据应用架构师}”素养，而非人人去从零编写单机数据库内核。这些核心素养首要包括数据建模与高效查询构建能力：即能够面对现实业务中纷繁杂乱、非结构化的原始数据集合，运用现代 DBMS 建立严谨规范的数据模式，并通过编程式接口与高级查询语言实现亚秒级的低延迟检索与一致性更新。

此外，现代计算机专业毕业生必须深入研读以下两大新兴维度：
\begin{itemize}
    \item \textbf{数据在现代组织与社会治理中的全生命周期角色}：
    \begin{itemize}
        \item \textbf{数据生命周期}：数据生成与采集$\to$ 处理与清洗$\to$ 审查与报表呈现$\to$ 留存、归档与检索$\to$ 销毁与销户清理。
        \item \textbf{数据资产的社会、伦理与法律维度}：数据爆炸带来的规模效应、数据隐私权益、内生合规与隐私设计、去标识化与匿名化技术、数据所有权与数据资产边界、数据可靠性与防篡改审计、全栈数据库安全，以及算法应用中预期之外的次生歧视或伦理偏见。
    \end{itemize}
    \item \textbf{超越传统单机关系系统的现代存储与大数据分析前沿技术}：
    \begin{itemize}
        \item 正在显著增强、补充甚至在特定高并发场景下替代传统关系系统的 NoSQL 数据库（如 JSON/XML 文档存储、高性能键值存储、图数据库等）；
        \item 云原生弹性分布式数据库；
        \item MapReduce 分布式海量批处理范式；
        \item 面向现代数据科学与分析工程的数据帧与自动化数据治理管道。
    \end{itemize}
\end{itemize}

\textbf{数据管理生态中的多元专业工程角色：}\\
大纲充分认识到参与现代数据管理链条的多元化分工角色：
\begin{itemize}
    \item \textbf{产品功能工程师}：在业务系统一线同时驾驭 SQL 关系数据库与 NoSQL 操作型数据库，支撑高并发核心业务功能。
    \item \textbf{数据分析工程师 / 数据工程师}：精通分析型 SQL、Python 和 Scala 编程，构建高质量数据仓库、数据湖资产与端到端 ETL/ELT 自动化数据管道。
    \item \textbf{业务分析师}：熟练使用电子表格、敏捷商业智能（BI）工具与低代码报表进行业务数据多维洞察。
    \item \textbf{数据基础架构工程师}：专注研发与维护高可靠、高吞吐的联机事务处理（OLTP）引擎、分布式存储基座与云原生查询加速器。
    \item \textbf{“全员共识” (Everyone)}：每一位生产或消费数据的计算专业人员，都必须深刻理解数据背后的法律、伦理、安全与社会责任。
\end{itemize}

\textbf{超越一切分工的核心担当——“数据托管人”：}\\
过去，数据仅仅被视为企业内部可供调配与增值的“普通资产”（信息系统管理视角）。而在今天，数据被赋予了不可剥夺的人权与法律属性：关于客户的数据本质上属于客户个体自身（或由客户保留所有权）。现代计算专业人员必须树立清醒共识：合乎法律伦理地存储、防护和使用机构数据，是从业者作为一名负责任的“数据托管人”的基本底线。

\textbf{科学理论探索与工程实践导向之间的张力调和：}\\
大纲深刻指出了计算机课程体系中“职业工程技能培养”与“作为科学理论严密探索”之间的内在张力。这一矛盾在数据管理领域尤为尖锐。例如，形式化证明函数依赖理论中阿姆斯特朗公理系统（Armstrong's Axioms）的完备性与可靠性，是经典数据库理论的学术瑰宝；然而，绝大多数毕业生在漫长的工程研发职业生涯中可能永远不会直接书写此类数学证明。类似的情况也存在于关系规范化高阶理论之中。反之，如果毕业生只会机械地运行自动化建模工具将模式切分至 BCNF 范式，只会编写低效嵌套的初级 SQL，却对底层物理索引在加速查询执行中的决定性作用一无所知，这无疑是计算机高等教育的重大失职。

为此，CS2023 在学时规划上推陈出新：将全专业通用的计算机核心学时精简为 9 学时，而将大量底层进阶机制划入领域核心（20.5 学时）。这种结构赋予了各高校灵活定制的自适应空间：对于偏向系统底层的培养方案，可深入钻研 OLTP 存储引擎、事务可串行化与故障恢复机制；而对于偏向应用开发与数据科学的方案，则可侧重既有大型分布式数据库的编程式访问、现代数据建模与高级查询优化。

\textbf{计算机高等教育的历史反讽：}\\
长期以来，计算机专业学生在本科期间的大部分时间里，都在学习通过平面文本文件读写或控制台交互式 I/O 来处理数据；然而在象牙塔外的真实产业世界中，绝大多数真实数据（以绝对压倒性的体量）都是通过应用程序编程接口（API、数据库驱动与 ORM）从数据库中访问与持久化的。CS2023 强烈呼吁高校应在程序设计早期尽早引入编程式数据库访问，并在整个大学课程体系中持续强化这一现代工程习惯。

\textbf{与网络安全 (SEC) 及社会伦理 (SEP) 的深度交织与原点地位：}\\
尽管数据管理在学科分类上与网络安全（SEC）及社会、伦理与专业素养（SEP）正交，但现实中它恰恰是安全与伦理的“风暴原点”。在设计持久化存储系统时，“什么数据应当被采集和存储”必须首先经受严格的法律合规与伦理道德审查：是否存在隐私侵蚀？数据安全防护体系是否固若金汤？
\end{preamblebox}

\section{核心学时分配建议}

本知识领域在“向日葵模型”下的核心学时推荐如下表所示：

\begin{table}[H]
\centering
\setlength{\tabcolsep}{3.5pt}\fontsize{8pt}{9.6pt}\selectfont\begin{tabularx}{\columnwidth}{Y c c}
\toprule
\textbf{知识单元} & \textbf{计算机核心} & \textbf{领域核心} \\
\midrule
数据的作用与治理 & 2 & -- \\
数据库系统核心概念 & 2 & 0.5 \\
数据建模 & 2 & 3 \\
关系数据库与规范化理论 & 1 & 3 \\
查询构建与 SQL 语言 & 2 & 4 \\
查询处理与执行优化 & -- & 4 \\
数据库管理系统内核实现 & -- & 4 \\
NoSQL 系统与非关系型存储 & -- & 2 \\
数据安全与隐私保护 & 待定 & 待定 \\
数据分析与大数据处理 & 待定 & 待定 \\
分布式数据库与云计算 & \emph{拓展选修} & \emph{拓展选修} \\
半结构化与非结构化数据库 & \emph{拓展选修} & \emph{拓展选修} \\
\midrule
\textbf{合计} & \textbf{9 学时} & \textbf{20.5 学时$^*$} \\
\bottomrule
\end{tabularx}
\begin{flushleft}\footnotesize
\vspace{-2pt}
$^*$注：大纲文本总计标记为 25 学时，其中包含为制定中的“数据安全与隐私”及“数据分析”模块预留的 4.5 个核心学时。
\end{flushleft}
\end{table}

\section{知识单元与详细大纲}

\subsection{数据的作用与治理}
\textbf{【学时建议】}：2 计算机核心学时。

本单元系统阐述数据在现代社会与企业组织中的战略资产地位，全面梳理数据的完整生命周期，并深入剖析数据采集、留存与使用过程中的法律、伦理与合规边界。

\subsubsection{计算机核心主题}
\begin{itemize}
    \item \textbf{数据生命周期}：
    \begin{itemize}
        \item 数据生成与采集：业务产生、传感器采集、日志埋点与跨系统汇聚。
        \item 处理与清洗：数据验证、异常值清洗、格式转换与丰富增强。
        \item 审查与报表呈现：商业智能展现、审计追踪与数据质量校验。
        \item 留存、归档与检索：冷热数据分级存储、合规归档与高效索引检索。
        \item 销毁与销户清理：安全删除机制、加密擦除、被遗忘权与不可逆物理擦除。
    \end{itemize}
    \item \textbf{数据集合的社会与法律伦理维度}：
    \begin{itemize}
        \item 数据规模化风险：海量数据集中带来的攻击面扩大与系统性风险。
        \item 数据隐私权益：个人身份信息（PII）保护、隐私泄漏防范。
        \item 内生合规与隐私设计：在数据库模式设计之初内嵌合规约束（如 GDPR、HIPAA 等）。
        \item 匿名化与去标识化：泛化、置换、差分扰动及抗再标识攻击能力。
        \item 数据所有权与资产边界：用户数据所有权、平台托管权与数据资产变现合法性。
        \item 数据可靠性与可信度：数据真实性、防篡改机制、数据溯源。
        \item 预期用途与非预期次生风险：数据二次利用、算法画像歧视与监控风险防范。
    \end{itemize}
\end{itemize}

\subsubsection{示范性学习成果}
\begin{enumerate}
    \item \textbf{[解释]} 阐明数据生命周期各核心阶段的划分逻辑，并列举确保各阶段数据完整性与合规性的工程措施。
    \item \textbf{[评估]} 针对给定的实际业务应用，全面评估其在海量数据采集与长期持久化留存中的隐私、法律与伦理风险，撰写合规防范白皮书。
\end{enumerate}

\subsection{数据库系统核心概念}
\textbf{【学时建议】}：2 计算机核心学时，0.5 领域核心学时。

本单元奠定数据库系统最本质的理论概念与体系结构认知，剖析数据库相对于传统操作系统的文件系统在并发、持久性与逻辑抽象层级上的革命性跨越。

\subsubsection{计算机核心主题}
\begin{itemize}
    \item \textbf{数据库系统的用途与核心优势}：消除数据冗余、保障数据一致性、提供原子并发控制、强故障恢复能力与声明式数据独立性。
    \item \textbf{数据库系统总体架构与组件划分}：存储管理器（文件组织、缓冲池管理）、查询处理器（解析器、优化器、执行引擎）、事务管理器（并发控制、恢复管理）。
    \item \textbf{DBMS 核心功能模块设计原理}：查询执行机制、事务一致性维护、缓冲池置换策略、底层物理访问路径。
    \item \textbf{三级模式体系结构与数据独立性}：外模式（视图层）、模式（逻辑层）与内模式（物理层）；物理数据独立性（物理存储变化不影响逻辑模式）与逻辑数据独立性（逻辑模式改变不破坏外部视图）。
    \item \textbf{声明式查询语言的核心思想}：关注“要获取什么数据”（What）而非“如何计算与获取数据”（How）。
    \item \textbf{事务管理基本概念}：ACID 特性（原子性 Atomicity、一致性 Consistency、隔离性 Isolation、持久性 Durability）的工程意义。
    \item \textbf{规范化理论初步}：消除重复与冗余数据的动因，更新异常、插入异常与删除异常的本质。
    \item \textbf{大规模海量数据管理策略}：水平扩展理念、NoSQL 系统破局思路、MapReduce 分布式并行处理范式。
    \item \textbf{针对纯增删改查业务应用的支持模式}：轻量级接口封装与高效对象映射。
    \item \textbf{分布式与云原生数据库系统概览}：弹性计算与存储分离架构、跨机房复制基础。
    \item \textbf{结构化、半结构化与非结构化数据库的分类与特征对比}。
\end{itemize}

\subsubsection{领域核心主题}
\begin{itemize}
    \item \textbf{流数据与现代流式数据库系统}：支持持续结构化流式内容摄入、滑动窗口聚合与亚秒级响应的流数据库架构。
\end{itemize}

\subsubsection{示范性学习成果}
\begin{enumerate}
    \item \textbf{[解释]} 辨析并准确陈述使用专业数据库系统相较于操作系统平面文件系统所提供的至少四项决定性优势。
    \item \textbf{[解释]} 列举关系数据库系统的核心架构组件，并清晰描述各组件的职责边界。
    \item \textbf{[应用]} 完整追踪一条典型 SQL 查询在数据库管理系统各组件之间解析、校验、优化、生成物理计划与检索磁盘的执行全过程。
    \item \textbf{[评估]} 论证数据独立性（尤其是物理数据独立性）对现代大型软件系统长期解耦演进的核心价值。
    \item \textbf{[开发]} 在 SQL 中熟练编写基础的多表“选择-投影-连接 (SPJ)”查询。
    \item \textbf{[解释]} 详细阐明确保事务正确执行的核心 ACID 四大特性及其违背时可能产生的破坏性后果。
    \item \textbf{[解释]} 阐述消除数据库中重复冗余数据的深层优势（存储优化、更新开销降低与数据不一致规避）。
    \item \textbf{[解释]} 概述 MapReduce 如何利用“数据不动计算动”与数据局部性原理实现海量数据的高效并行处理。
    \item \textbf{[评估]} 权衡对比结构化、半结构化与非结构化数据在数据模式灵活性、查询表达能力与存储效率上的本质差异。
\end{enumerate}

\subsection{数据建模}
\textbf{【学时建议】}：2 计算机核心学时，3 领域核心学时。

数据建模是将现实世界复杂的业务实体及其相互联系抽象转化为计算机可精确理解与操作的逻辑结构的关键桥梁。

\subsubsection{计算机核心主题}
\begin{itemize}
    \item \textbf{数据建模基础理论}：概念模型、逻辑模型与物理模型的划分与映射转换过程。
    \item \textbf{关系数据模型}：关系（表）、元组（行/记录）、属性（列/字段）、域（取值范围）、关系模式与关系实例的形式化定义。
\end{itemize}

\subsubsection{领域核心主题}
\begin{itemize}
    \item \textbf{概念数据模型构建}：
    \begin{itemize}
        \item 实体-联系模型（E-R 模型）：实体集、弱实体集、属性（单值/多值/复合/派生）、联系集及其度数（一元/二元/多元）。
        \item 映射基数与参与约束：$1:1$、$1:n$、$n:m$ 关联关系以及完全参与/部分参与约束。
        \item 统一建模语言（UML）类图在概念数据建模中的表达与转换。
    \end{itemize}
    \item \textbf{半结构化数据模型}：利用模式定义语言（如 DTD、XML Schema 或 JSON Schema）对半结构化文档进行严格的类型约束与层级结构规范化定义。
\end{itemize}

\subsubsection{拓展选修主题}
\begin{itemize}
    \item \textbf{电子表格数据模型}：二维网格表格计算模型与公式依赖拓扑。
    \item \textbf{面向对象数据模型}：对象标识符（OID）、封装性、继承体系与对象关系映射（ORM，关联 PL 领域）；GraphQL 统一图模式定义与图查询。
    \item \textbf{国际 SQL 标准演进前沿}：跟踪 SQL:2023 及后续标准最新引入的多维数组、属性图查询（SQL/PGQ）等新特性。
    \item \textbf{专用领域数据建模专题}：
    \begin{itemize}
        \item 时序数据建模：时间戳对齐、时序聚合、窗口多路连接与时间衰减归档。
        \item 图数据建模：节点属性、边属性、异构图拓扑建模与路径遍历。
        \item 物化视图与高维近似数据结构：物化视图维护策略；专用概率统计数据结构（如用于超大海量基数估计的 HyperLogLog、位图 Bitmap 等）。
        \item 性能架构警示：深入剖析为何直接对海量“原始高频时序数据”执行日均价、日去重 UV 统计、日活跃度等查询极易导致严重性能雪崩，掌握预聚合与分层汇总模式。
        \item 地理空间数据建模：点、线、多边形拓扑表示，空间参照系统与地理信息空间索引建模。
    \end{itemize}
\end{itemize}

\subsubsection{示范性学习成果}
\begin{enumerate}
    \item \textbf{[解释]} 准确阐述关系数据模型的核心构成要素（关系、元组、属性、域）及其代数本质。
    \item \textbf{[开发]} 熟练运用关系模型，通过主外键设计与中间关联表，准确建模现实业务中的 $1:1$、$1:n$ 和 $n:m$ 联系。
    \item \textbf{[解释]} 阐明实体-联系（E-R）模型及其他非关系概念模型的核心语义构件与基数约束表达机制。
    \item \textbf{[开发]} 面对给定的真实复杂业务环境，抽象出精确的概念数据模型（绘制标准 E-R 图或 UML 类图）。
    \item \textbf{[开发]} 针对高灵活性或松散结构的业务场景，构建规范的基于文档存储或键值存储的数据模式（如编写 JSON Schema）。
\end{enumerate}

\subsection{关系数据库与规范化理论}
\textbf{【学时建议】}：1 计算机核心学时，3 领域核心学时。

关系数据库以严密的数学集合论和一阶谓词逻辑为支撑。本单元聚焦关系完整性约束、模式映射以及运用函数依赖进行规范化设计的核心理论。

\subsubsection{计算机核心主题}
\begin{itemize}
    \item \textbf{关系完整性约束规则}：
    \begin{itemize}
        \item 实体完整性：主键非空且唯一。
        \item 参照完整性：外键引用有效性、级联操作。
        \item 键的形式化层次：超键、候选键与主键的定义与判定。
    \end{itemize}
    \item \textbf{关系数据库逻辑设计基础}：关系模式初步建立与约束声明。
\end{itemize}

\subsubsection{领域核心主题}
\begin{itemize}
    \item \textbf{概念模式向关系模式的映射转换}：强实体集、弱实体集、二元关联、多元联系、ISA 继承层次的标准化表转换算法。
    \item \textbf{物理数据库设计}：底层文件结构（堆文件、顺序文件、定长与变长记录）、块存储组织与记录指针映射。
    \item \textbf{函数依赖理论导引}：函数依赖形式化定义（$X \to Y$）、逻辑蕴涵与平凡/非平凡依赖。
    \item \textbf{关系规范化理论}：
    \begin{itemize}
        \item 关系模式分解准则：无损连接分解的判定定理；保持函数依赖的形式化检验。
        \item 核心范式体系：第一范式（1NF，属性不可分）、第二范式（2NF，消除非主属性对码的部分依赖）、第三范式（3NF，消除非主属性对码的传递依赖）与 Boyce-Codd 范式（BCNF，消除任何属性对码的部分与传递依赖）。
        \item 工业反规范化设计：针对高频读性能瓶颈，审慎引入受控冗余表、宽表与冗余字段的工程权衡与同步一致性开销。
    \end{itemize}
\end{itemize}

\subsubsection{拓展选修主题}
\begin{itemize}
    \item \textbf{高级函数依赖理论}：
    \begin{itemize}
        \item 阿姆斯特朗公理系统（自反律、增广律、传递律及其派生规则）的严密数学证明。
        \item 属性集闭包（$X^+$）的高效线性时间计算算法及其在超码判定中的应用。
        \item 最小函数依赖集 / 正规覆盖的极小化构造算法。
    \end{itemize}
    \item \textbf{高阶规范化理论与多值依赖}：
    \begin{itemize}
        \item 多值依赖 (MVD)的形式化定义与第四范式（4NF）。
        \item 连接依赖 (JD)与投射-连接范式（PJNF / 第五范式 5NF）。
        \item 关系数据库模式表示理论。
    \end{itemize}
\end{itemize}

\subsubsection{示范性学习成果}
\begin{enumerate}
    \item \textbf{[解释]} 深入阐明关系数据模型背后的本质数学特征（关系即元组集合的笛卡尔积子集）。
    \item \textbf{[解释]} 深入剖析外键、候选键与超键在唯一性保障和多表关联上的本质差异。
    \item \textbf{[解释]} 全面列举关系数据库所支持的完整性约束类型（域约束、主键约束、参照完整性、用户自定义 CHECK 约束）。
    \item \textbf{[开发]} 将包含 $1:1$、$1:n$ 和 $n:m$ 联系的概念数据模式严谨转换为无歧义的关系模式集合。
    \item \textbf{[应用]} 针对给定的关系表与物理索引，配置匹配的底层文件存储组织形式（如堆文件或聚簇文件）。
    \item \textbf{[解释]} 阐明函数依赖理论如何将“超键与候选键”的概念进行严谨的形式化推广。
    \item \textbf{[评估]} 运用代数算法（如表格判定法）严格证明或否定一个模式分解是否具备无损连接性与保持依赖性。
    \item \textbf{[评估]} 检测给定关系模式所属的最高范式等级（1NF、2NF、3NF 还是 BCNF），并指出潜在的数据更新异常。
    \item \textbf{[评估]} 深入论述在特定高吞吐工程场景下对关系表实施反规范化宽表设计的合理动机、收益及其数据同步风险。
\end{enumerate}

\subsection{查询构建与 SQL 语言}
\textbf{【学时建议】}：2 计算机核心学时，4 领域核心学时。

SQL 是与关系数据库进行人机和程序化交互的国际通用语言。本单元涵盖关系代数形式化基础、复杂查询表达、数据操纵与数据库编程式访问。

\subsubsection{计算机核心主题}
\begin{itemize}
    \item \textbf{SQL 声明式查询语句规范与构建}：
    \begin{itemize}
        \item 单表与多表查询：SELECT 子句、FROM 子句、WHERE 谓词过滤。
        \item 分组与聚合分析：GROUP BY、HAVING 过滤谓词、聚集函数（COUNT、SUM、AVG、MIN、MAX）。
        \item 多表连接技术：内连接、左/右/全外连接、自连接。
        \item 集合运算：UNION、INTERSECT、EXCEPT/MINUS。
        \item 嵌套子查询：标量子查询、集合成员测试、量化比较（ALL/ANY）、存在性测试。
    \end{itemize}
\end{itemize}

\subsubsection{领域核心主题}
\begin{itemize}
    \item \textbf{关系代数理论基础}：
    \begin{itemize}
        \item 经典集合运算：并（$\cup$）、交（$\cap$）、差（$-$）、广义笛卡尔积（$\times$）。
        \item 专门关系运算：选择（$\sigma$）、投影（$\pi$）、更名（$\rho$）、自然连接（$\bowtie$）、$\theta$-连接、条件等值连接、除法运算（$\div$）。
        \item 关系代数表达式的形式化推导与等价变换。
    \end{itemize}
    \item \textbf{SQL 全体系语言规范}：
    \begin{itemize}
        \item 数据定义语言（DDL）：CREATE/ALTER/DROP TABLE，主外键约束、唯一约束（UNIQUE）、CHECK 检查条件、非空约束与默认值。
        \item 数据更新语言（DML）：INSERT 插入、UPDATE 批量更新、DELETE 条件删除与 MERGE 操作；处理事务控制与隐式提交。
        \item 视图管理：视图定义、更新条件判定与可更新视图约束。
    \end{itemize}
\end{itemize}

\subsubsection{拓展选修主题}
\begin{itemize}
    \item \textbf{关系演算}：元组关系演算（TRC）与域关系演算（DRC）；安全关系演算表达式；Codd 等价定理（关系代数与安全关系演算的表达能力等价性）。
    \item \textbf{图形化与第四代查询环境}：示例查询 (QBE)与敏捷数据构建环境。
    \item \textbf{宿主语言调用非过程化查询的多种途径}：嵌入式 SQL（Embedded SQL）、动态 SQL、CLI/ODBC/JDBC 驱动标准、现代对象关系映射（ORM）框架的利弊权衡（关联 PL 领域）。
    \item \textbf{其他主流专用查询语言导引}：半结构化 XML 查询语言、语义网与资源描述框架图查询语言（SPARQL）、图数据库属性查询语言。
    \item \textbf{数据库端过程化扩展编程}：存储过程、用户自定义函数（UDF）、行级/语句级触发器及其级联调试。
\end{itemize}

\subsubsection{示范性学习成果}
\begin{enumerate}
    \item \textbf{[开发]} 熟练编写涵盖多表连接、分组聚合、集合并/交/差以及除法语义的高级 SQL 查询语句。
    \item \textbf{[评估]} 准确辨别一个嵌套子查询是否属于“相关子查询”，并评估其潜在的重复计算性能开销。
    \item \textbf{[应用]} 在通用宿主编程语言（如 Python、Java 等）中，通过标准数据库驱动连接池编程式执行参数化 SQL 查询，并高效迭代处理返回的元组结果集。
    \item \textbf{[开发]} 纯熟运用 SQL DDL 语法完整定义关系模式，涵盖所有主外键依赖、CHECK 业务约束以及级联更新/删除触发器。
    \item \textbf{[开发]} 编写精准、具备事务原子性保障的目标元组更新与安全批量变更 SQL 脚本。
\end{enumerate}

\subsection{查询处理与执行优化}
\textbf{【学时建议】}：4 领域核心学时。

数据库的高性能源自其底层复杂精密的物理存储索引、查询操作符算法以及基于代价的查询优化器（CBO）。本单元揭示查询引擎“黑盒”之下的物理运转机理。

\subsubsection{领域核心主题}
\begin{itemize}
    \item \textbf{核心物理索引结构}：
    \begin{itemize}
        \item B+ 树索引结构：节点结构（分支节点与叶子节点）、高扇出、自平衡分裂与合并算法、范围扫描的高性能保障。
        \item 哈希索引：静态哈希、动态可扩展哈希与线性哈希；点查优势与无法支持范围查询的物理局限。
        \item SQL 索引生命周期管理：普通索引、唯一索引、复合联合索引的最左前缀原则。
    \end{itemize}
    \item \textbf{查询操作符的物理实现算法}：
    \begin{itemize}
        \item 外部归并排序：两路/多路归并、内存缓冲区页帧划分、运行段生成与合并 I/O 代价。
        \item 选择操作的物理实现：全表顺序扫描vs 基于聚簇/非聚簇索引的高效点查与范围检索。
        \item 投影操作的物理实现：基于排序与基于哈希的重复记录消除机制。
        \item 自然连接与等值连接算法族：简单嵌套循环连接、块嵌套循环连接、排序归并连接、Grace 哈希连接。
        \item 算法复杂度与效率量化分析：以磁盘 I/O 块传输次数为度量标准的精确代价估算模型。
    \end{itemize}
    \item \textbf{查询代数等价变换规则}：选择操作尽早下推、投影操作提前消除冗余属性、连接顺序重排等启发式优化规则。
    \item \textbf{基于代价的查询优化器}：
    \begin{itemize}
        \item 物理访问路径选择。
        \item 物理执行计划生成：左深连接树与动态规划搜索空间剪枝（System R 算法）。
        \item 统计信息与选择率估算：直方图、数据基数与均匀性独立假设。
        \item 仅索引扫描计划（Index-only Plans / 覆盖索引 Covering Index）：免除回表物理 I/O 开销的终极优化方案。
    \end{itemize}
    \item \textbf{数据库系统性能调优与索引选型}：
    \begin{itemize}
        \item 索引对系统性能的双刃剑效应：加速只读查询吞吐 vs 显著拖慢高频插入/更新/删除操作并占用磁盘空间的权衡。
    \end{itemize}
\end{itemize}

\subsubsection{示范性学习成果}
\begin{enumerate}
    \item \textbf{[解释]} 阐明 B+ 树与动态哈希索引的内部拓扑组织结构，解释为什么关系数据库普遍采用 B+ 树作为默认主索引。
    \item \textbf{[开发]} 熟练编写 SQL DDL 语句创建高选择性的聚簇索引与复合索引，优化特定业务查询负载。
    \item \textbf{[解释]} 详细步进式剖析外部归并排序、去重投影、排序归并连接、散列哈希连接以及块嵌套循环连接的完整运行流程。
    \item \textbf{[评估]} 针对上述查询物理算法，根据分配的可用内存页数与关系表数据块规模，精确推导并计算其总磁盘 I/O 请求代价。
    \item \textbf{[应用]} 运用关系代数等价变换公理重写复杂 SQL 逻辑代数树，使其转化为适合流水线执行的物理执行计划。
    \item \textbf{[评估]} 根据字段统计直方图与谓词条件，精确计算单表过滤及两表等值连接的选择率估计值。
    \item \textbf{[开发]} 阐述如何通过在复合索引中追加包含字段，使高频查询能够直接在索引树上完成所有字段投影（覆盖索引扫描）。
    \item \textbf{[评估]} 针对复杂现实业务场景下的读写负载特征，权衡并决策出支持全套高频查询所需保留的最精炼索引组合。
\end{enumerate}

\subsection{数据库管理系统内核实现}
\textbf{【学时建议】}：4 领域核心学时。

DBMS 内核是系统软件设计的集大成者。本单元剖析内存缓冲池置换、并发事务的隔离性与调度可串行化理论，以及保障持久性与原子性的故障恢复子系统。

\subsubsection{领域核心主题}
\begin{itemize}
    \item \textbf{数据库缓冲池管理}：
    \begin{itemize}
        \item 内存缓冲区页帧组织、脏页标记与固定计数器。
        \item 经典置换策略：最近最少使用（LRU）、时钟算法、LRU-K 以及扫描抗击防护机制。
        \item 脏页异步后台刷盘调度与预读优化。
    \end{itemize}
    \item \textbf{事务处理与并发语义}：
    \begin{itemize}
        \item 事务隔离级别全解析：读未提交、读已提交、可重复读与可串行化。
        \item 并发读写异常本质：脏读、不可重复读与幻读。
        \item 事务 ACID 四大特性的底层技术支撑映射。
        \item 可串行化理论：调度、冲突等价、冲突可串行化判定准则与拓扑排序优先图。
    \end{itemize}
    \item \textbf{并发控制机制}：
    \begin{itemize}
        \item 封锁技术：排他锁（X 锁）与共享锁（S 锁）、意向锁（意向共享 IS、意向排他 IX）。
        \item 两阶段封锁协议（2-Phase Locking, 2PL）：生长阶段与收缩阶段，严格两阶段锁（Strict 2PL）防范级联回滚，强两阶段锁（Rigorous 2PL）。
        \item 死锁应对策略：死锁预防（Wait-Die 与 Wound-Wait 时间戳算法）、死锁检测（等待图 Wait-For Graph 周期环检测）与超时回滚。
    \end{itemize}
    \item \textbf{恢复管理器与缓冲池交互}：
    \begin{itemize}
        \item 缓冲池刷盘管理策略：STEAL/NO-STEAL（是否允许未提交事务将脏页刷盘）与 FORCE/NO-FORCE（是否要求提交事务在返回前将所有脏页完全落盘）。
    \end{itemize}
\end{itemize}

\subsubsection{拓展选修主题}
\begin{itemize}
    \item \textbf{高级并发控制协议}：
    \begin{itemize}
        \item 乐观并发控制 (OCC)：读取阶段、有效性确认验证阶段与写入阶段。
        \item 时间戳排序协议 (T/O)与多版本并发控制（MVCC，快照隔离 Snapshot Isolation 的底层版本链与 GC）。
    \end{itemize}
    \item \textbf{故障恢复子系统深度实现}：
    \begin{itemize}
        \item 预写日志技术 (WAL)：日志先行核心原则，日志序列号（LSN，PageLSN 与 FlushedLSN 判定）。
        \item 工业级 ARIES 故障恢复算法体系：
        \begin{itemize}
            \item 分析阶段：从检查点重构活跃事务表（ATT）与脏页表（DPT）。
            \item 重做阶段：从最早未落盘脏页开始重放历史，恢复系统至崩溃瞬间状态。
            \item 撤销阶段：反向回滚未提交事务，生成补偿日志记录（CLR）防止恢复期间再次崩溃。
        \end{itemize}
    \end{itemize}
\end{itemize}

\subsubsection{示范性学习成果}
\begin{enumerate}
    \item \textbf{[解释]} 阐明 DBMS 缓冲池管理器的内存分页管理机制、缓冲置换算法与页固定机制。
    \item \textbf{[解释]} 深入阐释一个正确的事务管理器所必须提供的 ACID 核心特性，以及数据库内核如何通过并发控制与 WAL 日志保障这四项特性。
    \item \textbf{[解释]} 概述可串行化的数学理论基石，利用事务读写冲突规则构建优先图，准确判定一个调度是否具备冲突可串行化。
\end{enumerate}

\subsection{系统与非关系型存储}
\textbf{【学时建议】}：2 领域核心学时。

互联网高并发、海量弹性扩展及半结构化数据的爆发催生了 NoSQL 运动。本单元探讨关系模型以外的替代性数据组织形态及其在分布式环境下的设计权衡。

\subsubsection{领域核心主题}
\begin{itemize}
    \item \textbf{为何需要 NoSQL 系统？}：
    \begin{itemize}
        \item 传统关系模型与现代应用程序对象模型之间的“阻抗失配”。
        \item 传统单机 RDBMS 在面对超大规模海量读写时的横向水平扩展物理瓶颈。
        \item 动态无模式或松散模式对敏捷业务迭代的强力支撑。
    \end{itemize}
    \item \textbf{两大主流非关系型数据模型}：
    \begin{itemize}
        \item 键值数据模型：极简的 PUT/GET/DELETE 原语、极高读写吞吐与内存化加速（如 Redis）。
        \item 文档数据模型：以 JSON/BSON 为核心组织单元的层次嵌套半结构化表达（如 MongoDB）。
    \end{itemize}
\end{itemize}

\subsubsection{拓展选修主题}
\begin{itemize}
    \item \textbf{非关系存储底层引擎实现}：基于 LSM 树（Log-Structured Merge-tree，MemTable、SSTable、Compaction）的高吞吐写优化存储引擎。
    \item \textbf{分布式分布模型}：
    \begin{itemize}
        \item 数据分片：范围分片、哈希分片与一致性哈希路由。
        \item 数据副本复制：单主多从同步/异步复制、无主对等复制。
    \end{itemize}
    \item \textbf{分布式一致性模型与理论极限}：
    \begin{itemize}
        \item CAP 定理：一致性、可用性、分区容错性的三者权衡边界。
        \item BASE 最终一致性模型：基本可用、软状态、最终一致。
        \item Quorum 一致性读写仲裁协议（$N, W, R$ 参数配置与强/弱一致性切换）。
    \end{itemize}
    \item \textbf{分布式并行处理范式}：MapReduce 分布式计算范式、多阶段与增量 MapReduce。
    \item \textbf{工业案例研究与工程反模式}：
    \begin{itemize}
        \item 云原生对象存储系统（如 AWS S3 等海量非结构化资产底座）。
        \item 图数据库系统（Graph Databases，复杂社交拓扑与风控网络）。
        \item 深度工程反思：“\textbf{何时绝不应使用 NoSQL}”——识别缺乏 ACID 强一致性、涉及多表强关联金融级结算场景中的 NoSQL 滥用反模式。
    \end{itemize}
\end{itemize}

\subsubsection{示范性学习成果}
\begin{enumerate}
    \item \textbf{[解释]} 结合真实的工程业务场景，清晰论证在特定需求下选择使用 NoSQL 替代传统 RDBMS 的决定性用例与权衡动因。
    \item \textbf{[解释]} 深入阐释键值存储与文档数据库背后的本质特征、数据组织拓扑结构及其在读写吞吐上的核心优势。
\end{enumerate}

\subsection{数据安全与隐私保护}
\textbf{【学时建议】}：学时待定（CS2023 框架正与 SEC 知识领域深度统筹规划）。

数据管理是现代网络空间安全与个人隐私保护的最核心阵地。本单元涵盖从物理静态存储到逻辑访问层的全链路安全加固。

\subsubsection{大纲规划重点议题}
\begin{itemize}
    \item \textbf{数据加密与完整性审计}：静态透明数据加密 (TDE)、传输层 TLS/SSL 加密、密钥生命周期管理与列级敏感字段端到端加密。
    \item \textbf{细粒度访问控制体系}：基于角色的访问控制（RBAC）、基于属性的动态访问控制（ABAC）、数据库行级安全 (RLS)与动态数据脱敏。
    \item \textbf{数据库隐私保护计算前沿}：差分隐私（Differential Privacy）在统计聚合查询中的拉普拉斯/高斯噪声注入机制；多方安全计算（SMPC）与同态加密在分布式联合查询中的应用。
    \item \textbf{数据匿名化工程标准}：$k$-匿名（$k$-Anonymity）、$l$-多样性（$l$-Diversity）与 $t$-贴近度（$t$-Closeness）在防止准标识符重识别攻击中的实现。
    \item \textbf{数据库抗攻击与审计防线}：SQL 注入漏洞的编程式根除（强参数化查询与预编译语句）、数据库操作审计日志与防篡改区块链存证。
\end{itemize}

\subsection{数据分析与大数据处理}
\textbf{【学时建议】}：学时待定（CS2023 框架正与 ACM 数据科学课程体系全面对齐）。

从联机事务处理（OLTP）走向联机分析处理（OLAP）与智能决策支持，是现代数据技术演进的另一条主线。

\subsubsection{大纲规划重点议题}
\begin{itemize}
    \item \textbf{数据仓库与联机分析处理 (OLAP)}：星型模型、雪花模型与事实表/维度表设计；多维数据集与上卷、下钻、切片和切块。
    \item \textbf{面向海量分析的物理存储革命}：列式存储格式（Columnar Storage，如 Parquet、ORC）、列式压缩算法（字典编码、Run-Length 游程编码、位图编码）与向量化查询执行引擎。
    \item \textbf{现代数据工程治理与管道}：自动化 ETL/ELT 管道架构、数据湖与湖仓一体架构；数据血缘追踪与数据质量度量。
    \item \textbf{现代计算抽象与分布式框架}：基于 DataFrame 的声明式数据转换、分布式内存并行分析引擎（如 Apache Spark）。
\end{itemize}

\subsection{分布式数据库与云计算}
\textbf{【学时建议】}：拓展选修。

面对全球跨地域容灾与海量伸缩需求，现代数据库已经全面拥抱分布式与云原生架构。

\subsubsection{拓展选修主题}
\begin{itemize}
    \item \textbf{分布式数据库系统}：
    \begin{itemize}
        \item 分布式数据存储：水平分段、垂直分段、混合分段与全局目录管理。
        \item 分布式查询处理与优化：分布式查询重写、半连接算法与最小化网络传输开销策略。
        \item 分布式事务原子性协议：两阶段提交协议 (2PC)的协调者/参与者状态转移、单点阻塞故障与三阶段提交（3PC）演进。
        \item 同构与异构分布式数据库集成：联邦数据库系统与跨异构源中间件。
        \item 客户端-服务器分布式计算拓扑（跨领域关联 系统基础/计算范式）。
    \end{itemize}
    \item \textbf{并行数据库系统架构}：
    \begin{itemize}
        \item 经典物理架构分类：共享内存、共享磁盘与无共享架构的伸缩性与容错性对比。
        \item 性能扩展度量：加速比与扩展比的理论上限。
        \item 数据多副本高可用同步与弱一致性容灾模型。
    \end{itemize}
\end{itemize}

\subsection{半结构化与非结构化数据库}
\textbf{【学时建议】}：拓展选修。

现代人工智能与多媒体应用的繁荣带来了对高维向量、复杂嵌套文档与空间数据的爆发式存储检索需求。

\subsubsection{拓展选修主题}
\begin{itemize}
    \item \textbf{高维向量化非结构化数据与向量存储}：
    \begin{itemize}
        \item 非结构化文本与多媒体特征提取：基于 N-gram 的 TF-IDF 文本特征向量化、Word2Vec 词向量嵌入。
        \item 专用向量数据库与多维数组数据库：支持高维向量近似最近邻搜索（ANN，如 HNSW、IVF-PQ 算法）。
    \end{itemize}
    \item \textbf{半结构化文档数据深度管理（以 JSON 为核心）}：
    \begin{itemize}
        \item 物理存储引擎：嵌套半结构化数据类型的高效二进制编码与列存压缩（如 JSONB、BSON）。
        \item 专用索引加速技术：
        \begin{itemize}
            \item 跳跃索引与布隆过滤器加速块级过滤。
            \item 倒排索引与位图压缩索引。
            \item 面向半结构化地理空间数据的空间填充曲线索引（Space Filling Curve，如 Z-order 曲线、Hilbert 空间填充曲线）。
        \end{itemize}
        \item 混合负载（HTAP）下的查询处理权衡：插入吞吐、范围点查、就地更新/删除之间的工程权衡。
        \item 经典前沿系统工程剖析：PostgreSQL 的 JSONB 扩展支持、MongoDB 纯文档架构与 Snowflake 云数仓半结构化 JSON 混合查询。
    \end{itemize}
\end{itemize}

\section{专业素养与思维倾向}

数据管理不仅是一门关于算法与存储工程的技术学科，更直接承载着现代数字社会的信任基石。本领域着重塑造学生以下四项关键专业思维特质：
\begin{itemize}
    \item \textbf{专业严谨}：这是数据管理中最为核心的职业素养。尽管全球不同国家与法域对“数据所有权”的法律界定存在差异，但任何现代企业与计算专业人员对于严密、合规、专业化管理数据的敬畏之心与职业操守绝不可动摇。
    \item \textbf{责任担当}：在专业处理海量（个人敏感）数据时，承担高度社会责任至关重要。从业者必须具备两重底线：既要构筑坚如磐石的技术防线抵御未经授权的非法访问与网络攻击，更要从伦理高度坚决杜绝“看似合法合规但本质上不负责任、侵蚀公众利益的数据滥用行为”。此外，数据托管人必须确保数据在面对意外崩溃、物理灾难与不可抗力风险时能够完好无损地复原。
    \item \textbf{协作精神}：数据库管理员、架构师、安全审计员与业务数据消费者之间必须建立高度透明、无缝对接的协作机制，确保在合规授权的前提下让正确的数据以合乎规范的方式流通并赋能业务。
    \item \textbf{需求导向与敏捷响应}：数据的存储模式、保留周期与索引结构选型绝非闭门造车，而必须紧密敏捷地响应组织真实的业务需求与合规监管变化。
\end{itemize}

\section{数学基础要求}

为深入理解与推导关系代数、函数依赖与范式分解理论，学生必须牢固掌握以下核心数学前置知识：
\begin{itemize}
    \item \textbf{离散朴素集合论}：
    \begin{itemize}
        \item 集合的基本概念、元素属于与子集包含关系；
        \item 集合代数运算：并集、交集、差集；
        \item 广义笛卡尔积与元组映射；
        \item 幂集概念及其在属性集闭包计算中的应用。
    \end{itemize}
\end{itemize}

\section{交叉主题与跨学科共享概念}

数据管理领域与计算机科学各大知识领域存在着深厚而频繁的交叉共振：
\begin{itemize}
    \item \textbf{社会、伦理与专业素养 (SEP)}：数据隐私立法、GDPR/CCPA 时代合规审计、被遗忘权与算法公平性。
    \item \textbf{程序设计语言 (PL/FPL)}：声明式语法解析、类型系统、对象关系阻抗失配、ORM 映射与内存对象持久化。
    \item \textbf{系统基础 (SF)}：现代计算机系统的存储层次结构（Cache、RAM、NVMe SSD、HDD）、页式内存管理与 I/O 调度。
    \item \textbf{算法与复杂性 (AL)}：B+ 树平衡检索、多路归并外部排序、哈希冲突消解与动态规划查询优化。
    \item \textbf{并行与分布式计算 (PDC)}：数据分片路由、多副本复制协议、分布式事务提交（2PC）、分布式锁与 MapReduce 并行批处理。
    \item \textbf{专用平台开发 (SPD)}：移动端嵌入式数据库、物联网边缘时序存储。
    \item \textbf{计算机安全 (SEC)}：静态存储加密（TDE）、细粒度行级访问控制（RLS）、SQL 注入攻击防范与审计溯源。
    \item \textbf{数据挖掘与数据科学}：海量特征工程与湖仓一体分析（系统参见 ACM 数据科学核心规范）。
\end{itemize}

\section{课程体系规划与开课建议}

针对我国高校不同的办学定位与专业特色，建议采取以下三种典型的开课与课程规划路径：
\begin{enumerate}
    \item \textbf{标准计算机专业必修单门模式（“数据库系统原理与设计”，4 学分 / 64 学时）}：
    全面覆盖全部 9 个计算机核心学时以及 20.5 个领域核心内容。理论部分涵盖概念数据建模、关系代数与 SQL、函数依赖与 BCNF/3NF 规范化理论、B+ 树与哈希索引、基于代价的查询优化、事务并发控制（2PL）与 WAL 恢复机制；实验部分必须包含真实数据库应用系统的端到端设计开发，并包含针对复杂查询的索引性能调优分析。
    \item \textbf{高阶双轨递进模式（“基础数据库系统” $+$ “海量数据管理与分布式存储”）}：
    适用于高水平研究型大学与卓越工程师培养计划。第一门课专注经典关系理论与 SQL 开发应用；第二门进阶选修课系统探讨 NoSQL 存储引擎、分布式一致性协议、分布式事务（2PC）、列式存储与现代云原生数据仓库。
    \item \textbf{先行贯通融合模式}：
    破除传统“前两年完全用纯文本读写文件，直到大三才接触数据库”的脱节弊端。强烈建议在大学第一年《程序设计基础》（SDF）课程中尽早引入轻量级嵌入式数据库（如 SQLite）的编程式访问，培养学生使用标准 SQL 进行数据持久化存储的现代工程素养。
\end{enumerate}

\section{真实工程能力任务规范}

本领域制定了 8 个涵盖系统开发、风险审计、架构设计、性能调优与故障恢复的典型真实能力任务：

\begin{taskbox}[任务 1：为业务目标实现数据库编程式访问与数据处理]
\textbf{【任务情境】}：现代软件系统需要通过应用程序代码连接数据库，执行自动化数据检索与业务计算，而非依赖人工交互式控制台操作。\\
\textbf{【能力陈述】}：能够通过标准编程式接口（如 JDBC/ODBC、ORM 或原生数据库驱动）安全连接数据库，构造参数化查询，并在宿主编程语言中遍历返回的关系元组流，执行相应的数据处理与业务计算逻辑。\\
\textbf{【能力领域】}：软件工程、系统基础 \quad \textbf{【能力单元】}：系统开发、测试验证\\
\textbf{【关联合约/前置知识】}：数据管理/查询构建、社会伦理/社会语境、社会伦理/公平多元与包容、社会伦理/伦理分析方法、社会伦理/专业素养与工程伦理\\
\textbf{【核心要求】}：计算机核心 \quad \textbf{【达标等级】}：\textbf{开发}
\end{taskbox}

\begin{taskbox}[任务 2：评估新老应用数据收集与留存的社会及伦理风险（白皮书撰写）]
\textbf{【任务情境】}：组织拟上线一款涉及大规模用户敏感数据的新系统（或升级既有业务系统），需在立项之初评估其法律合规与伦理影响。\\
\textbf{【能力陈述】}：撰写专业决策白皮书，全面识别多方利益相关者，深入评估数据采集与长期持久化留存的长期潜在后果；审慎分析非预期数据滥用及潜在数据泄露事件所带来的危害与法律风险，并提出合规缓解方案。\\
\textbf{【能力领域】}：系统管理、理论基础 \quad \textbf{【能力单元】}：系统评估、数据治理、社会议题调适\\
\textbf{【关联合约/前置知识】}：社会伦理/社会语境、社会伦理/伦理分析方法、社会伦理/隐私与公民自由、社会伦理/专业伦理、社会伦理/安全法规与计算机犯罪、社会伦理/公平多元与包容、数据管理/数据的作用、安全隐私/安全基础\\
\textbf{【核心要求】}：计算机核心 \quad \textbf{【达标等级】}：\textbf{评估与解释}
\end{taskbox}

\begin{taskbox}[任务 3：面向新应用场景评估并决策底层数据存储架构模型]
\textbf{【任务情境】}：面对新型业务形态（如高并发社交互动、物联网大规模时序流水或金融级强事务支付），需在关系模型与非关系模型间做出科学架构选型。\\
\textbf{【能力陈述】}：综合剖析应用的数据读写特征（读写比、ACID 强一致性需求 vs BASE 最终一致性、水平扩展性诉求与模式可变性），在传统关系数据库、键值存储、文档数据库、图数据库之间进行定量权衡，选定最适合的数据模型并提供详尽的技术论证。\\
\textbf{【能力领域】}：应用架构、理论基础 \quad \textbf{【能力单元】}：系统评估、架构设计\\
\textbf{【关联合约/前置知识】}：数据管理/数据的作用、数据管理/数据库系统核心概念、数据管理/数据建模、数据管理/关系数据库、数据管理/NoSQL 系统、社会伦理/安全法规与计算机犯罪\\
\textbf{【核心要求】}：计算机核心与领域核心 (计算机核心 / 领域核心) \quad \textbf{【达标等级】}：\textbf{评估与解释}
\end{taskbox}

\begin{taskbox}[任务 4：基于索引调优与查询计划优化数据库应用执行性能]
\textbf{【任务情境】}：企业核心业务数据库在流量高峰期响应缓慢，关键业务报表与核心高频交易查询发生严重 I/O 阻塞。\\
\textbf{【能力陈述】}：分析数据库执行计划与查询物理访问路径，评估既有索引的利用率与更新维护代价；精准创建适宜的高效索引结构（如复合 B+ 树索引、动态哈希索引、覆盖索引/仅索引扫描）或移除低效冗余索引，显著降低高频重要查询的 I/O 次数与响应延迟。\\
\textbf{【能力领域】}：软件工程、应用开发 \quad \textbf{【能力单元】}：系统设计、系统重构与性能改进\\
\textbf{【关联合约/前置知识】}：数据管理/查询处理、数据管理/DBMS 内核实现、并行分布式/并行算法分析与编程\\
\textbf{【核心要求】}：领域核心 \quad \textbf{【达标等级】}：\textbf{评估与开发}
\end{taskbox}

\begin{taskbox}[任务 5：基于视图与访问控制机制实现数据安全防护与权限隔离]
\textbf{【任务情境】}：跨部门业务系统需共享底层数据库表，但必须严格防范越权访问与未经授权的敏感隐私字段泄露。\\
\textbf{【能力陈述】}：根据最小特权原则与数据治理合规要求，设计并创建精细化数据库视图，结合行级/列级访问控制策略和用户权限角色划分，从数据访问层彻底杜绝非授权访问。\\
\textbf{【能力领域】}：系统管理 \quad \textbf{【能力单元】}：运维管理、系统管理\\
\textbf{【关联合约/前置知识】}：数据管理/数据的作用、数据管理/关系数据库、数据管理/查询处理、社会伦理/安全法规与计算机犯罪、社会伦理/专业伦理、社会伦理/隐私与公民自由、安全隐私/安全基础\\
\textbf{【核心要求】}：领域核心 \quad \textbf{【达标等级】}：\textbf{开发}
\end{taskbox}

\begin{taskbox}[任务 6：在系统灾难或突发中断（如断电/崩溃）后快速恢复数据库至安全状态]
\textbf{【任务情境】}：生产数据库服务器突发断电、操作系统宕机或硬件故障，导致内存缓冲池未刷盘事务中断，需保障数据一致性并迅速重新上线。\\
\textbf{【能力陈述】}：掌握数据库恢复管理器工作机制，基于预写日志（WAL）与检查点技术，熟练运行恢复与还原工具，执行事务的 REDO（重做）与 UNDO（撤销）流程，将数据库迅速恢复到故障前最新的持久化一致安全状态。\\
\textbf{【能力领域】}：系统基础 \quad \textbf{【能力单元】}：运维管理、系统管理\\
\textbf{【关联合约/前置知识】}：数据管理/DBMS 内核实现\\
\textbf{【核心要求】}：领域核心 \quad \textbf{【达标等级】}：\textbf{应用与实践}
\end{taskbox}

\begin{taskbox}[任务 7：面向新业务应用完成从概念建模到关系规范化存储设计]
\textbf{【任务情境】}：组织启动一项新业务应用系统的研发，需要从零构建结构清晰、无数据异常的高质量底层数据存储方案。\\
\textbf{【能力陈述】}：深入理解业务领域需求，运用实体-联系模型（E-R 图）或 UML 类图完成概念数据建模；将概念模式规范映射为关系模式（准确定义主键、候选键、外键与完整性约束），并应用函数依赖与规范化理论将其分解至 BCNF/3NF 范式，消除插入、删除与更新异常。\\
\textbf{【能力领域】}：应用架构 \quad \textbf{【能力单元】}：系统设计、模式改进\\
\textbf{【关联合约/前置知识】}：数据管理/数据的作用、数据管理/数据建模、社会伦理/社会语境、社会伦理/伦理分析方法、社会伦理/隐私与公民自由、社会伦理/专业伦理、社会伦理/安全法规与计算机犯罪、社会伦理/公平多元与包容、安全隐私/安全基础\\
\textbf{【核心要求】}：计算机核心与领域核心 (计算机核心 / 领域核心) \quad \textbf{【达标等级】}：\textbf{应用与开发}
\end{taskbox}

\begin{taskbox}[任务 8：端到端完整实现新应用数据库系统的工程构建与部署落地]
\textbf{【任务情境】}：新系统从立项到落地需要一体化贯通数据存储需求分析、合规治理、模型设计、系统部署与索引优化的全生命周期。\\
\textbf{【能力陈述】}：综合执行端到端数据库工程构建全流程：完成概念与逻辑数据建模；系统审计数据采集与存储的伦理法律风险；科学决断关系型与 NoSQL 存储选型；编写全套 DDL 脚本定义数据表、完整性约束与触发器；针对核心查询负载合理配置 B+ 树与哈希索引，完成整套生产级数据库实例的构建与验证。\\
\textbf{【能力领域】}：架构设计、软件工程、系统基础 \quad \textbf{【能力单元】}：系统开发、测试验证\\
\textbf{【关联合约/前置知识】}：数据管理/数据的作用、数据管理/数据库系统核心概念、数据管理/数据建模、数据管理/关系数据库、数据管理/NoSQL 系统、数据管理/DBMS 内核实现、SEP 各单元、安全隐私/安全基础\\
\textbf{【核心要求】}：计算机核心与领域核心 (计算机核心 / 领域核心) \quad \textbf{【达标等级】}：\textbf{开发}
\end{taskbox}

\section{指导委员会名单}

\textbf{委员会主席}：
\begin{itemize}
    \item Mikey Goldweber, 泽维尔大学
\end{itemize}

\textbf{委员会委员}：
\begin{itemize}
    \item Sherif Aly, 开罗美国大学
    \item Sara More, 约翰霍普金斯大学
    \item Mohamed Mokbel, 明尼苏达大学
    \item Rajendra Raj, 罗切斯特理工学院
    \item Avi Silberschatz, 耶鲁大学 —— \emph{注：国际知名计算机科学家，经典著作《操作系统概念》与《数据库系统概念》资深作者}
    \item Min Wei, 微软公司
    \item 向乔，厦门大学
\end{itemize}

\begin{strip}
\section{附录：核心主题与技能要求等级对照表}

\noindent\begin{tabularx}{\textwidth}{p{0.18\textwidth} Y c c c}
\toprule
\textbf{知识单元} & \textbf{核心知识点与具体主题} & \textbf{认知等级} & \textbf{分类} & \textbf{学时} \\
\midrule
\multirow{2}{*}{\makecell[l]{数据管理/数据的作用\\与治理}} 
& 数据生命周期（创建、处理、审查、归档、销毁）与全链路合规 & 评估 & 计算机核心 & \multirow{2}{*}{2} \\
& 数据集合的社会与法律伦理维度（规模效应、隐私设计、所有权） & 评估 & 计算机核心 & \\
\midrule
\multirow{3}{*}{\makecell[l]{数据管理/数据库系统\\核心概念}} 
& 数据库系统核心价值与架构组件；三级模式结构与数据独立性 & 解释 & 计算机核心 & \multirow{2}{*}{2} \\
& 事务管理、规范化消除冗余、海量数据 MapReduce 批处理机制 & 解释 & 计算机核心 & \\
& 声明式查询语言使用范式与基础 SPJ 复杂查询构建 & 开发 & 计算机核心 & \\
\cline{2-5}
& 现代结构化流数据与流式计算数据库系统 & 解释 & 领域核心 & 0.5 \\
\midrule
\multirow{2}{*}{\makecell[l]{数据管理/数据建模}} 
& 现实世界数据建模流程；关系数据模型（关系、元组、属性） & 开发 & 计算机核心 & 2 \\
\cline{2-5}
& 概念数据建模（E-R 图、UML 类图）；半结构化模型 & 解释 & 领域核心 & 3 \\
\midrule
\multirow{2}{*}{\makecell[l]{数据管理/关系数据库与\\规范化理论}} 
& 实体完整性、参照完整性（外键约束）；超键与候选键定义 & 解释 & 计算机核心 & 1 \\
\cline{2-5}
& 概念模式向关系模式转换；底层物理文件与存储结构设计；函数依赖理论；规范化范式（BCNF/3NF 无损连接与保持依赖分解）；反规范化设计 & 开发 & 领域核心 & 3 \\
\midrule
\multirow{2}{*}{\makecell[l]{数据管理/查询构建与\\SQL 语言}} 
& SQL 声明式查询高级构建（单表/多表过滤、分组聚合、集合操作） & 开发 & 计算机核心 & 2 \\
\cline{2-5}
& 关系代数（并、交、差、笛卡尔积、选择、投影、连接、除法）；SQL DDL 模式完整性定义；SQL DML 更新子语言 & 开发 & 领域核心 & 4 \\
\midrule
\multirow{2}{*}{\makecell[l]{数据管理/查询处理与\\执行优化}} 
& 物理索引结构（B+ 树、哈希索引）；查询操作符算法（外部排序、选择、去重投影、连接算法族）；等价代数变换；基于代价的优化器（访问路径、选择率估算、仅索引扫描） & 解释 & 领域核心 & \multirow{2}{*}{4} \\
& 数据库物理性能调优：工作负载分析与索引选型决策（权衡读写开销） & 开发 & 领域核心 & \\
\midrule
\multirow{2}{*}{\makecell[l]{数据管理/DBMS\\内核实现}} 
& 缓冲池管理（页置换、固定机制）；事务处理与 ACID 特性保障 & 解释 & 领域核心 & \multirow{2}{*}{4} \\
& 事务隔离级别与可串行化理论；两阶段封锁协议（2PL）与死锁处理；恢复管理器与缓冲池交互机制 & 解释 & 领域核心 & \\
\midrule
\multirow{2}{*}{\makecell[l]{数据管理/NoSQL 系统与\\非关系型存储}} 
& 为何需要 NoSQL？CRUD 应用与关系数据库的阻抗失配；水平扩展动因 & 解释 & 领域核心 & \multirow{2}{*}{2} \\
& 键值数据模型与基于 JSON 的文档数据模型核心特征与适用边界 & 解释 & 领域核心 & \\
\midrule
数据管理/数据安全与隐私 & 规划中：静态加密、细粒度行级访问控制、差分隐私、脱敏防泄露 & 制定中 & 规划中 & 待定 \\
\midrule
数据管理/数据分析 & 规划中：OLAP 架构、列式存储、数据湖仓一体与现代数据管道 & 制定中 & 规划中 & 待定 \\
\bottomrule
\end{tabularx}
\end{strip}
